% Input note supplied by the PI on 30 September 2026. Incorporated through
% docs/UNIFIED_MODEL_RETRAINING_PLAN_2026-09-30.md ("Incorporating the supplied
% outlier note"), which records what was adopted (empirical contaminant
% background, calibrated QSO-support percentile) and what was not (the binary
% QSO-vs-contaminant score as ranking statistic). Kept for provenance; not a
% list of open tasks. Body unchanged since 30 September.
\documentclass[11pt]{article}

\usepackage[margin=1in]{geometry}
\usepackage{amsmath}
\usepackage{amssymb}
\usepackage{bm}
\usepackage{hyperref}
\usepackage[round,authoryear]{natbib}
\usepackage{enumitem}

\title{Robust Probabilistic Quasar Selection with Extreme Deconvolution:\\
Empirical Contaminant Modelling and Protection Against Gaussian-Tail Extrapolation}
\author{}
\date{}

\begin{document}
\maketitle

\begin{abstract}
We consider probabilistic quasar selection using extreme-deconvolution (XD)
models of quasars and contaminants in colour or relative-flux space. A
failure mode arises when an object lies far from both the stellar and quasar
loci. Because Gaussian mixtures have non-zero support everywhere, both class
likelihoods can be extremely small while their ratio is nevertheless very
large. Such an object can therefore receive a high nominal quasar probability
even though it is not plausibly drawn from the quasar population. This is
fundamentally a closed-set classification problem rather than a numerical
pathology of XD. We review how the XD/XDQSO literature, Bayesian quasar
selection, empirical contaminant modelling, and the broader
reject-option/open-set literature address the underlying issue. We recommend
that our classifier replace the narrow ``QSO versus star'' comparison by
``QSO versus all contaminants'', where the contaminant density is learned
empirically from a representative sample of the survey catalogue. This should
be combined with a separately calibrated QSO-support or out-of-distribution
criterion. The posterior then answers whether an object is more plausibly a
QSO than an object normally encountered in the survey, while the support
statistic prevents remote Gaussian tails from being interpreted as positive
evidence for QSO membership.
\end{abstract}

\section{The problem}

Let $\bm{x}$ denote the true colour or relative-flux vector of an object,
$\bm{y}$ its measured value, and $\bm{S}$ the measurement-error covariance
matrix. For a class $C$, an XD model represents the intrinsic density as a
Gaussian mixture,
%
\begin{equation}
p_C(\bm{x}\mid m)
=
\sum_{k=1}^{K_C}
\alpha_{Ck}(m)
\mathcal{N}
\left(
\bm{x}\mid
\bm{\mu}_{Ck}(m),
\bm{V}_{Ck}(m)
\right),
\end{equation}
%
where $m$ denotes a reference magnitude, or more generally a variable on
which the population density is conditioned. XD was introduced specifically
to infer such latent distributions in the presence of heteroscedastic
measurement uncertainties and missing data \citep{bovy2011xd}.

For an object with measurement covariance $\bm{S}$, the predictive likelihood
under the Gaussian-error approximation is
%
\begin{equation}
L_C(\bm{y},\bm{S}\mid m)
=
\sum_{k=1}^{K_C}
\alpha_{Ck}(m)
\mathcal{N}
\left(
\bm{y}\mid
\bm{\mu}_{Ck}(m),
\bm{V}_{Ck}(m)+\bm{S}
\right).
\label{eq:xdlike}
\end{equation}

Suppose that the classifier contains only quasars, $Q$, and stars, $S$. A
natural posterior is then
%
\begin{equation}
P(Q\mid\bm{y},\bm{S},m)
=
\frac{
n_Q(m)L_Q(\bm{y},\bm{S}\mid m)
}{
n_Q(m)L_Q(\bm{y},\bm{S}\mid m)
+
n_S(m)L_S(\bm{y},\bm{S}\mid m)
},
\label{eq:closedposterior}
\end{equation}
%
where $n_Q(m)$ and $n_S(m)$ encode the population number densities or prior
weights.

Equation~(\ref{eq:closedposterior}) contains the failure mode of interest. The
posterior depends on the \emph{relative} weighted likelihoods, not on whether
either likelihood is appreciable in absolute terms. For example, let
%
\begin{equation}
n_Q L_Q = \epsilon,
\qquad
n_S L_S = \epsilon^2 .
\end{equation}
%
Then
%
\begin{equation}
P(Q\mid\bm{y})
=
\frac{1}{1+\epsilon}
\longrightarrow 1
\end{equation}
%
as $\epsilon\rightarrow 0$, although
%
\begin{equation}
n_Q L_Q+n_S L_S\longrightarrow 0 .
\end{equation}

Thus an object can receive a QSO posterior arbitrarily close to unity while
being extremely improbable under both fitted populations. This is not an
error in Bayes' theorem. It is a consequence of the model assumption that the
classes appearing in the denominator exhaust the possible populations.

The effect is particularly easy to encounter with Gaussian mixtures because
a Gaussian has infinite support. An XD mixture therefore extrapolates a
non-zero density arbitrarily far from the training locus. If the stellar
mixture happens to decline faster than the QSO mixture in a particular
direction, the likelihood ratio can become very large even though the QSO
density itself is negligible.

\section{What the established literature implies}

\subsection{Extreme deconvolution and XDQSO}

Extreme deconvolution is a density-estimation method, not an open-set
classifier \citep{bovy2011xd}. Its role is to infer a latent population
density while correctly accounting for heterogeneous measurement errors.

The direct astronomical precedent is XDQSO
\citep{bovy2011xdqso}. XDQSO models the distributions of quasars and
contaminants in flux or relative-flux space, reconvolving the inferred latent
density with the uncertainty of each test object. It also combines the
density model with realistic population abundances rather than using colour
density alone.

The important conceptual point for the present problem is that a Bayesian
posterior has its intended interpretation only when the denominator contains
the relevant alternatives. A two-class $Q$--$S$ posterior is really a
probability conditional on the hypothesis that the source is in either the
QSO population or the population represented by the stellar model. If the
catalogue also contains compact galaxies, unusual stars, blends, moving
objects, problematic photometry, peculiar spectral energy distributions, or
other populations that are not represented by either model, the closed-set
interpretation is wrong.

A second lesson from XDQSO is that class priors matter. The relevant quantity
for target selection is closer to a number density in observable space than
to a unit-normalized colour density. Schematically,
%
\begin{equation}
q_C(\bm{y},\bm{S},m)
=
n_C(m)\,L_C(\bm{y},\bm{S}\mid m),
\end{equation}
%
and posterior odds should compare these weighted densities.

\subsection{Bayesian rare-object selection}

The Bayesian high-redshift QSO selection framework of
\citet{mortlock2012} makes the same general point from a different direction.
For a rare population, a candidate should be assessed by combining a model
for the observables with the expected abundance of the target and contaminant
populations. Quasar-like colours alone are not sufficient because contaminants
can vastly outnumber the target population.

This reinforces two requirements for our problem:

\begin{enumerate}[label=(\arabic*)]
\item the denominator should represent the populations that actually occur in
the input catalogue; and
\item the different populations should carry realistic prior number densities.
\end{enumerate}

\subsection{The most direct precedent: an empirical all-inclusive contaminant model}

The strongest direct precedent for the solution required here is the XDHZQSO
method of \citet{nanni2022}. They use XD to estimate the density of
high-redshift QSOs and the density of contaminants in multi-band flux space.
Crucially, the contaminants are not restricted to a narrowly defined stellar
class. The contaminant distribution is trained empirically from a large sample
of real survey sources, and the contaminating populations are grouped into a
single all-inclusive class.

This choice directly addresses the weakness of a QSO-versus-star classifier.
One does not need to know whether every false candidate is a cool dwarf,
compact galaxy, blend, or some other population. For classification, one needs
the denominator to describe the objects that the survey actually presents to
the target-selection algorithm.

Nanni et al. also emphasize the advantage of learning poorly understood
contaminants directly from the data, while retaining XD's principled treatment
of photometric uncertainties \citep{nanni2022}. This is very close to the
appropriate formulation for our application.

\subsection{Reject options, open-set recognition, and OOD detection}

There is a broader statistical context to the same problem. Classical work on
classification with a reject option formalizes the idea that a classifier
should sometimes abstain rather than force an assignment
\citep{chow1970}. In modern terminology, the presence of test objects outside
the set of training populations is an open-set recognition problem
\citep{scheirer2013}.

The connection to our case is straightforward. The current classifier is
closed-set: it assumes that all relevant objects are represented by its QSO
and stellar populations. Objects that violate that assumption can nevertheless
receive extreme posterior probabilities.

Generative likelihoods also require care when used for
out-of-distribution (OOD) detection. In modern generative-model work,
likelihood ratios against an explicit background model have been used to
reduce misleading likelihood behaviour on OOD data
\citep{ren2019}. The exact machine-learning setting differs from XDQSO, but
the conceptual lesson is relevant: relative likelihoods are most useful when
the comparison model represents a meaningful background population.

For our application, the literature therefore motivates two complementary
operations:

\begin{enumerate}[label=(\arabic*)]
\item replace the incomplete stellar denominator by a realistic empirical
contaminant model; and
\item retain a separate support/rejection test so that objects poorly
described by the population models are not forced into a known class.
\end{enumerate}

\section{Recommended primary classifier}

\subsection{QSO versus an empirical contaminant population}

Our preferred formulation is
%
\begin{equation}
P_Q
=
P(Q\mid\bm{y},\bm{S},m)
=
\frac{
n_Q(m)L_Q(\bm{y},\bm{S}\mid m)
}{
n_Q(m)L_Q(\bm{y},\bm{S}\mid m)
+
n_C(m)L_C(\bm{y},\bm{S}\mid m)
},
\label{eq:qcont}
\end{equation}
%
where $C$ is an \emph{all-inclusive empirical contaminant population}.

Operationally, $C$ should mean:

\begin{quote}
the distribution of non-QSO objects that can appear in the catalogue to
which the classifier will actually be applied.
\end{quote}

This is preferable to interpreting $C$ as a particular physical class. It may
contain ordinary stars, unusual stars, unresolved galaxies, compact galaxies,
blends, peculiar objects, and any other real source population that survives
the parent-sample cuts.

The main score should therefore answer
%
\begin{quote}
Is this object more plausibly drawn from the QSO population than from the
empirical population of non-QSO sources encountered by the survey?
\end{quote}
%
rather than merely
%
\begin{quote}
Is the extrapolated QSO Gaussian mixture larger than the extrapolated stellar
Gaussian mixture at this point?
\end{quote}

\subsection{Why not simply add a random ``other'' class on top of the stars?}

A representative random catalogue contains ordinary stars. If we keep the
existing stellar population and add an ``other'' model trained on an
unfiltered random survey sample, the stellar population is represented in both
terms and is therefore effectively double-counted.

There are two coherent alternatives.

The simpler and recommended option is to use a single all-inclusive
contaminant model $C$ and compare QSO versus contaminant as in
equation~(\ref{eq:qcont}).

If separate stellar probabilities are scientifically important, a genuine
three-class model can instead be written
%
\begin{equation}
P_Q
=
\frac{
n_Q L_Q
}{
n_Q L_Q+n_S L_S+n_O L_O
}.
\label{eq:threeclass}
\end{equation}
%
However, $S$ and $O$ must then represent distinct population components.
A raw random catalogue cannot simply be used as $O$ because it already
contains $S$.

A principled three-class implementation would require either labelled
partitioning or a semi-supervised mixture in which the QSO and stellar
components are constrained by their labelled training sets and an additional
flexible component absorbs the residual survey density. That is a more
complicated inference problem and is unnecessary if the scientific goal is
only robust QSO ranking.

\section{How to construct the empirical contaminant model}

\subsection{Match the deployment selection}

The contaminant training sample should be drawn from the same parent
population to which the classifier will be applied. In practice it should use
the same:

\begin{itemize}
\item detection requirements;
\item magnitude limits;
\item morphology requirements, if morphology is part of the parent selection;
\item required photometric coverage;
\item basic quality cuts; and
\item any broad pre-selection that defines the population entering the
classifier.
\end{itemize}

This point is essential. A generative posterior is only calibrated for the
population from which its training and prior models are derived. If the
contaminant sample is cleaner or otherwise differently selected than the
deployment catalogue, the model can recreate exactly the same missing-class
problem in a subtler form.

Catastrophic measurements that would never be considered usable targets can
be removed by reproducible quality cuts. However, it is generally a mistake
to clean the contaminant sample so aggressively that rare but real survey
objects disappear. Those sources are precisely part of the background against
which QSO candidates should be judged.

\subsection{Use a large representative sample}

After applying the parent-sample definition, draw a large random sample of
survey sources over the magnitude, depth, and sky ranges relevant for
deployment.

Known high-confidence QSOs should ideally be removed. Because QSOs are rare,
a small residual QSO fraction may have little effect on the fitted contaminant
density, but this should be checked rather than assumed.

If the catalogue properties vary strongly with field, depth, or Galactic
position, either the training set should reproduce those variations or the
model/prior should be conditioned on the variables responsible for them.

\subsection{Use the same XD noise treatment}

The QSO and contaminant models should be evaluated consistently. If the
intrinsic contaminant density is
%
\begin{equation}
p_C(\bm{x}\mid m)
=
\sum_{k=1}^{K_C}
\beta_k
\mathcal{N}
\left(
\bm{x}\mid
\bm{\nu}_k,
\bm{W}_k
\right),
\end{equation}
%
then candidate $i$ should have
%
\begin{equation}
L_{C,i}
=
\sum_{k=1}^{K_C}
\beta_k
\mathcal{N}
\left(
\bm{y}_i\mid
\bm{\nu}_k,
\bm{W}_k+\bm{S}_i
\right).
\end{equation}

The same measurement covariance convention must be used for the QSO model.
Otherwise the Bayes factor contains a mismatch in the treatment of
uncertainty.

It is often useful to condition or bin the density model by a reference
magnitude rather than make one Gaussian mixture fit the entire number-count
distribution. This follows the general XDQSO strategy
\citep{bovy2011xdqso}.

\subsection{Use realistic number-count priors}

The classifier should compare weighted population densities, not only
unit-normalized colour densities:
%
\begin{align}
q_Q &= n_Q(m,\bm{\theta}) L_Q(\bm{y},\bm{S}\mid m),\\
q_C &= n_C(m,\bm{\theta}) L_C(\bm{y},\bm{S}\mid m),
\end{align}
%
with
%
\begin{equation}
P_Q=\frac{q_Q}{q_Q+q_C}.
\end{equation}

Here $\bm{\theta}$ denotes additional variables that substantially affect the
population prior. Examples could include sky position, local source density,
or survey depth, if those effects are large enough to matter.

For contaminants, $n_C$ can usually be measured directly from the survey
after the parent-sample selection. The QSO prior can come from a
selection-corrected QSO training sample, a luminosity function, or another
population model appropriate for the science case.

A realistic low QSO prior is necessary, but it is not sufficient to solve the
tail problem. For any finite prior odds, a sufficiently extreme likelihood
ratio can still dominate. Population priors and support/OOD rejection address
different issues.

\section{Why the empirical contaminant model is not sufficient by itself}

Replacing ``star'' by ``all contaminants'' should eliminate most pathological
high scores, because objects that genuinely occur in unusual parts of survey
colour space will now be represented in the denominator.

It does not, however, eliminate the mathematical problem completely. Any
finite Gaussian mixture still has non-zero density everywhere. An object can
lie so far from both the QSO and empirical contaminant populations that both
likelihoods are negligible, while their ratio is again extreme.

The classifier should therefore be allowed to distinguish

\begin{quote}
QSO-like relative to the survey background
\end{quote}

from

\begin{quote}
actually supported by the QSO population.
\end{quote}

This motivates a second, independent support statistic.

\section{A practical QSO-support statistic}

For candidate $i$, define
%
\begin{equation}
\ell_{Q,i}
=
\log L_Q(\bm{y}_i,\bm{S}_i\mid m_i).
\end{equation}

A fixed numerical cut on $\ell_Q$ is unattractive because its scale depends on
dimension, population width, photometric covariance, and the chosen
coordinates. Instead, calibrate the QSO likelihood against the distribution
expected for genuine QSOs.

For fixed conditioning variables and measurement covariance, define
%
\begin{equation}
u_Q
=
{\rm Pr}_{\widetilde{\bm{y}}\sim Q}
\left[
\log L_Q(\widetilde{\bm{y}})
\leq
\log L_Q(\bm{y})
\right].
\label{eq:support}
\end{equation}

Thus $u_Q$ is the fraction of genuine-QSO realizations that would have lower
QSO predictive density than the candidate. Objects deep in an extrapolated
Gaussian tail have $u_Q$ close to zero.

A practical object-by-object Monte-Carlo procedure is:

\begin{enumerate}[label=(\arabic*)]
\item select or sample the appropriate intrinsic QSO mixture for the
candidate's magnitude or other conditioning variables;
\item draw latent QSO observables from that mixture;
\item add noise using the candidate covariance $\bm{S}_i$;
\item evaluate the same convolved QSO likelihood for every simulated
measurement; and
\item rank the observed $\ell_{Q,i}$ within this reference distribution.
\end{enumerate}

An independent held-out spectroscopic QSO sample can be used instead of, or
in addition to, simulations. Empirical validation is especially valuable
because a misspecified QSO mixture can make a formally calibrated simulation
look better than the real population.

The final selection should then use two distinct requirements:
%
\begin{equation}
P_Q > P_{\rm min}
\end{equation}
%
and
%
\begin{equation}
u_Q > u_{\rm min}.
\end{equation}

The first tests QSO versus survey contaminants. The second tests whether the
candidate lies in a region of observable space that is actually supported by
the QSO model.

The thresholds should be selected from held-out data according to the
scientific loss function: desired QSO completeness, available follow-up
density, purity, or some combination thereof. They should not be chosen from
the numerical magnitude of the raw Gaussian-mixture likelihood.

\section{Optional global model-support test}

A complementary quantity is the total fitted population evidence,
%
\begin{equation}
E_i
=
n_Q L_{Q,i}+n_C L_{C,i},
\end{equation}
%
or
%
\begin{equation}
\ell_{{\rm all},i}
=
\log\left(n_Q L_{Q,i}+n_C L_{C,i}\right).
\end{equation}

The distribution of this statistic can be calibrated on a held-out random
survey sample. Very low values indicate objects that are poorly described by
the entire fitted model rather than merely being poor QSO matches.

Such objects should be flagged as ``unmodelled'' or ``out of distribution''
instead of being forced into either QSO or contaminant. This global check is
useful for discovering catalogue artifacts, unusual populations, or failures
of the contaminant density model.

The QSO-support statistic $u_Q$ and the global-support statistic answer
different questions. For protecting QSO candidate selection against the
specific high-ratio/low-density failure mode, $u_Q$ is the more direct test.
The global support statistic is primarily a model-diagnostic and
abstention tool.

\section{Useful diagnostics}

Before defining the final selection, the following plots should be examined
on independent data.

\subsection{$\log L_Q$ versus $\log L_C$}

This is the most direct view of the generative comparison. The original
failure mode corresponds to objects for which both likelihoods are extremely
small. The old star-only classifier may have assigned them high QSO scores
because $\log L_S$ was even smaller.

\subsection{$P_Q$ versus $u_Q$}

This separates relative classification from absolute QSO support. The
pathological region is
%
\begin{equation}
P_Q\ {\rm high},
\qquad
u_Q\ {\rm low}.
\end{equation}

After a good empirical contaminant model is introduced, this population
should shrink substantially. The remaining objects are exactly those for
which an explicit support cut is valuable.

\subsection{Old score versus new score}

Compare the original QSO-versus-star posterior or log-odds with the new
QSO-versus-contaminant posterior. Objects whose old score was created mainly
by an extrapolated vanishing stellar likelihood should move strongly downward.

\subsection{Mixture-component responsibilities}

For QSO Gaussian $k$, define the responsibility
%
\begin{equation}
r_{ik}
=
\frac{
\alpha_k
\mathcal{N}
(\bm{y}_i\mid\bm{\mu}_k,\bm{V}_k+\bm{S}_i)
}{
L_{Q,i}
}.
\end{equation}

Inspecting $r_{ik}$ for pathological candidates can reveal whether a single
broad, low-weight Gaussian is responsible for much of the remote-tail
likelihood. This is a valuable model diagnostic, particularly for identifying
QSO-training contaminants or excessive mixture covariances.

It should not, however, replace the contaminant and support treatment. Every
finite Gaussian mixture has infinite support even when all components are
well fitted.

\section{Validation strategy}

\subsection{Held-out QSOs}

Use QSOs that were not fitted by the XD model to measure:

\begin{itemize}
\item the distribution of $P_Q$;
\item the distribution of $u_Q$;
\item completeness as a function of magnitude;
\item completeness as a function of redshift;
\item completeness as a function of photometric uncertainty; and
\item whether unusual but genuine parts of the QSO locus receive
systematically low support.
\end{itemize}

The last point is critical. A support criterion is only useful if the QSO
density model itself adequately represents the intrinsic diversity of QSOs.

\subsection{Held-out random survey objects}

Use a random survey sample that was not used to train the empirical
contaminant XD model. Run the full classifier on it without first selecting
objects by their QSO score.

This sample directly tests how frequently the deployed classifier generates
extreme false-positive scores in the real catalogue population.

In particular, inspect the sources satisfying
%
\begin{equation}
P_Q > P_{\rm min}
\end{equation}
%
before the support cut, and determine which of them have very small $u_Q$.
These are the objects closest to the original failure mode.

Where spectroscopy exists, purity and contamination can be measured directly.
Where it does not, the held-out survey sample still provides an unbiased
estimate of the frequency and morphology of extreme scores under the actual
deployment distribution.

\section{Approaches that should not be the primary solution}

\subsection{A stricter QSO/star posterior threshold}

Increasing the threshold on the original two-class posterior does not solve
the mathematical problem. If
%
\begin{equation}
L_Q/L_S\rightarrow\infty,
\end{equation}
%
the posterior can approach unity even when $L_Q\rightarrow 0$.

\subsection{A hard floor on the stellar likelihood}

A rule such as
%
\begin{equation}
L_S\rightarrow\max(L_S,L_{\rm floor})
\end{equation}
%
can numerically limit the likelihood ratio, but the floor has no clear
population interpretation. It effectively introduces an undocumented
background class.

If such a floor is required for numerical stability, it should not be
confused with the statistical solution. The scientifically interpretable
solution is to model the relevant contaminant/background population.

\subsection{Covariance shrinking of the QSO mixture}

Making the QSO Gaussian components narrower can reduce the frequency of
remote-tail candidates, particularly if an excessively broad component is
poorly fitted. However, covariance tuning does not solve the conceptual
problem because every Gaussian still has infinite support.

The mixture structure should be selected by predictive validation, not tuned
specifically to suppress objects that expose a closed-set assumption.

\subsection{QSO priors alone}

A realistic QSO prior is essential in rare-object selection
\citep{bovy2011xdqso,mortlock2012}. It does not guarantee protection against
a pathological likelihood ratio. A finite prior penalty can always be
overcome if the contaminant likelihood becomes sufficiently smaller than the
QSO likelihood.

\section{Recommended implementation for our classifier}

We recommend the following concrete workflow.

\begin{enumerate}[label=\arabic*.]

\item Retain the current XD QSO density model, subject to predictive validation
on QSOs not used in fitting.

\item Define precisely the parent catalogue to which the classifier is
applied.

\item Draw a large representative random sample from that catalogue using the
same parent-sample selection, magnitude range, photometric coverage, and
baseline quality requirements.

\item Remove secure known QSOs where practical.

\item Fit an XD contaminant density to this sample. The contaminant model
should use the same observables and measurement-error convention as the QSO
model.

\item Estimate the contaminant number counts $n_C$ directly from the survey
and adopt a justified QSO prior $n_Q$.

\item Replace the current QSO-versus-star score by
%
\begin{equation}
\boxed{
P_Q
=
\frac{n_Q L_Q}
{n_Q L_Q+n_C L_C}
}.
\end{equation}

\item Retain the old stellar likelihood as a diagnostic if useful, but do not
let the stellar model alone define the non-QSO denominator.

\item Using held-out QSOs and/or Monte-Carlo realizations, calibrate a
QSO-support percentile $u_Q$ that accounts for each object's measurement
covariance.

\item Require both a sufficiently large $P_Q$ and a sufficiently large $u_Q$.
Choose the thresholds from independent validation data to achieve the desired
science trade-off between completeness, purity, and target density.

\item Optionally calculate a global model-support statistic based on
$n_Q L_Q+n_C L_C$ and flag objects that are inconsistent with the full
population model.

\item Monitor the results with $\log L_Q$--$\log L_C$, $P_Q$--$u_Q$, and
old-score--new-score diagnostic plots.
\end{enumerate}

\section{Summary}

The high QSO scores obtained for sources far from the QSO locus are a
predictable consequence of closed-set generative classification. The statement
%
\begin{quote}
``the extrapolated QSO density is much larger than the extrapolated stellar
density''
\end{quote}
%
does not imply
%
\begin{quote}
``the source is plausibly drawn from the QSO population.''
\end{quote}

The most directly relevant astronomical precedent is to model the target
population against an all-inclusive empirical contaminant population, as done
in XDHZQSO \citep{nanni2022}, while using realistic population priors as in
probabilistic QSO-selection work more generally
\citep{bovy2011xdqso,mortlock2012}.

For our application, the primary classifier should therefore be
%
\begin{equation}
\boxed{
P_Q
=
\frac{n_Q L_Q}
{n_Q L_Q+n_C L_C}
}
\end{equation}
%
where $C$ is fitted from representative real survey sources.

Because Gaussian mixtures retain infinite tails, this should be accompanied
by a separate QSO-support requirement,
%
\begin{equation}
\boxed{
u_Q > u_{\rm min}
}.
\end{equation}

These two quantities answer different questions. $P_Q$ asks whether the
object is more consistent with QSOs than with the real survey background.
$u_Q$ asks whether the object is actually located in a region of observable
space supported by the QSO model. Using both is the principled way to prevent
objects in remote extrapolated tails from becoming high-priority QSO
candidates solely because the stellar or contaminant model falls even faster.

\begin{thebibliography}{99}

\bibitem[Bovy, Hogg \& Roweis(2011)]{bovy2011xd}
Bovy, J., Hogg, D. W., \& Roweis, S. T. 2011,
``Extreme deconvolution: Inferring complete distribution functions from
noisy, heterogeneous and incomplete observations,''
\emph{The Annals of Applied Statistics}, 5, 1657--1677.
\href{https://doi.org/10.1214/10-AOAS439}{doi:10.1214/10-AOAS439}.
Preprint: \url{https://arxiv.org/abs/0905.2979}.

\bibitem[Bovy et al.(2011)]{bovy2011xdqso}
Bovy, J., Hennawi, J. F., Hogg, D. W., Myers, A. D.,
Kirkpatrick, J. A., Schlegel, D. J., Ross, N. P., Sheldon, E. S.,
McGreer, I. D., Schneider, D. P., \& Weaver, B. A. 2011,
``Think Outside the Color Box: Probabilistic Target Selection and the
SDSS-XDQSO Quasar Targeting Catalog,''
\emph{The Astrophysical Journal}, 729, 141.
\href{https://doi.org/10.1088/0004-637X/729/2/141}
{doi:10.1088/0004-637X/729/2/141}.
Preprint: \url{https://arxiv.org/abs/1011.6392}.

\bibitem[Chow(1970)]{chow1970}
Chow, C. K. 1970,
``On Optimum Recognition Error and Reject Tradeoff,''
\emph{IEEE Transactions on Information Theory}, 16, 41--46.
\href{https://doi.org/10.1109/TIT.1970.1054406}
{doi:10.1109/TIT.1970.1054406}.

\bibitem[Mortlock et al.(2012)]{mortlock2012}
Mortlock, D. J., Patel, M., Warren, S. J., Hewett, P. C.,
Venemans, B. P., McMahon, R. G., \& Simpson, C. 2012,
``Probabilistic Selection of High-Redshift Quasars,''
\emph{Monthly Notices of the Royal Astronomical Society}, 419, 390--410.
\href{https://doi.org/10.1111/j.1365-2966.2011.19710.x}
{doi:10.1111/j.1365-2966.2011.19710.x}.
Preprint: \url{https://arxiv.org/abs/1101.4965}.

\bibitem[Nanni et al.(2022)]{nanni2022}
Nanni, R., Hennawi, J. F., Wang, F., Yang, J., Schindler, J.-T.,
\& Fan, X. 2022,
``Paving the way for Euclid and JWST via probabilistic selection of
high-redshift quasars,''
\emph{Monthly Notices of the Royal Astronomical Society},
515, 3224--3248.
\href{https://doi.org/10.1093/mnras/stac1944}
{doi:10.1093/mnras/stac1944}.

\bibitem[Ren et al.(2019)]{ren2019}
Ren, J., Liu, P. J., Fertig, E., Snoek, J., Poplin, R.,
DePristo, M., Dillon, J., \& Lakshminarayanan, B. 2019,
``Likelihood Ratios for Out-of-Distribution Detection,''
\emph{Advances in Neural Information Processing Systems}, 32.
\url{https://proceedings.neurips.cc/paper/9611-likelihood-ratios-for-out-of-distribution-detection}.
Preprint: \url{https://arxiv.org/abs/1906.02845}.

\bibitem[Scheirer et al.(2013)]{scheirer2013}
Scheirer, W. J., Rocha, A., Sapkota, A., \& Boult, T. E. 2013,
``Toward Open Set Recognition,''
\emph{IEEE Transactions on Pattern Analysis and Machine Intelligence},
35, 1757--1772.
\href{https://doi.org/10.1109/TPAMI.2012.256}
{doi:10.1109/TPAMI.2012.256}.

\end{thebibliography}

\end{document}
